\paragraph{Conjugate-coordinate class times a fuzzy factor: the tree non-derivative terms of the static scalar channel (conditional; bounded no-go and conditional route).}
This paragraph supplements the preceding one, in which the tree non-derivative terms of the wave were not computed. Keep its setting and notation, the tree action \((4g^2)^{-1}\int\rho_M\langle F_{ab}F^{ab}-6X_aX^a\rangle\) on \(M\), whose supplied bilinear term is that of the class with coefficient \(-3\) and not the mass term on \(K\), with \(C_{\mu\nu\lambda}=0\) there, the field equation \(\mathsf E_a=0\) and its linearisation \(L\) of the paragraph on the tree-level quadratic form of the level-one amplitudes, and a base point \(y_0\) with \(y^i=0\); write \(\mathsf P^a{}_b=y_0^ay_{0b}/(y_0\cdot y_0)\), \(\mathsf P_\perp=1-\mathsf P\), \(|y_0|^2=-y_0\cdot y_0\), and \(\langle\cdot\rangle\) for the fibre average.

\emph{Localisation.} For a wave \(A^a=\epsilon(y)\,\mathsf w^a\), with \(\mathsf w^a\) a fixed function on the leaf and \(\epsilon\) a function of the base functions of compact support, \(D^a\epsilon=s\,m\,\eta^{a\mu}\partial_\mu\epsilon\), and the symmetry of \(L\) under the Liouville pairing gives exactly
\[
\int\rho_M\langle A\cdot L(A)\rangle=\int\rho_M\Bigl[\epsilon^2\langle\mathsf w\cdot L(\mathsf w)\rangle+m^2\bigl(\eta^{\mu\nu}\partial_\mu\epsilon\,\partial_\nu\epsilon\,\langle\mathsf w\cdot\mathsf w\rangle-\langle(\mathsf w^\mu\partial_\mu\epsilon)^2\rangle\bigr)\Bigr],
\]
since \(2fL(fA)-f^2L(A)-L(f^2A)=2(Df\cdot Df)A-2Df\,(Df\cdot A)\) for every function \(f\) (exact in the polynomial Lie--Poisson algebra). The tree form of such a wave is therefore its leading symbol plus the single non-derivative term \(\epsilon^2\langle\mathsf w\cdot L(\mathsf w)\rangle\), with no term of first order in the derivatives of \(\epsilon\), and its first-order term \(\int\rho_M\langle\mathsf E\cdot A\rangle\) vanishes. For constant \(\kappa\), exactly and for both signs \(s\), \(L(\kappa T)=-2\kappa T+2(\operatorname{tr}\kappa)T\) for symmetric and \(L(\kappa T)=0\) for antisymmetric \(\kappa\); on the orbit \(T\cdot T=\ell^2+m^2=|y_0|^2\), and on the fibre over \(y_0\) the rotation invariance of the fibre measure gives \(\langle T^aT^b\rangle=\tfrac13(\ell^2+m^2)\mathsf P_\perp^{ab}\) (checked numerically).

\emph{The channel along \(y^\mu y^\nu\).} The channel has the realisation \(\mathsf w=\mathsf PT\), the rescaling \(X_{y_0}\to(1+\epsilon)X_{y_0}\) of the component \(X_{y_0}=y_{0a}X^a\). Its frame is \((1+\epsilon\mathsf P)E\), so that \(\delta\gamma^{\mu\nu}=2m^2\epsilon\,\mathsf P^{\mu\nu}\), \(t=2m^2\epsilon\), \(\sigma_1=-\epsilon\) and \(\sqrt{|g|}=m^{-4}(1+\epsilon)^{-1}\) exactly, the parametrisation of the preceding paragraph. Here \(L(\mathsf PT)=2\mathsf P_\perp T\), so \(\langle\mathsf w\cdot L(\mathsf w)\rangle=0\) on the whole leaf, and \(\mathsf w\) vanishes on the fibre over \(y_0\) since \(y\cdot T=0\), so the leading symbol vanishes there as well. For constant \(\epsilon\) this holds to all orders: with the matter fields at their vanishing background every term of the quartic term with \(C_{\mu\nu\lambda}=0\), of the supplied bilinear term and of the terms with the matrices of \(K\) is of degree two or zero in one component, since \([X_{y_0},X_{y_0}]=0\), so along the rescaling the action is \(S_\perp+(1+\epsilon)^2S_{y_0}\), also at finite matrix size, and stationarity gives \(S_{y_0}=0\); in the Poisson limit the change of the density is a total derivative for every \(\epsilon\) (exact polynomial identity; checked at finite size on fuzzy spheres, with the bilinear coefficients fixed by the field equation, for both contractions). Hence \(T_{\mathrm{tree}}=C_{\mathrm{tree}}=0\) in this realisation. That \(T_{\mathrm{tree}}=0\) holds in every realisation, since the background is stationary, so the tree terms do not cancel the tadpole of \(U_4\) in any channel; that \(C_{\mathrm{tree}}=0\) is a property of this realisation, on which the other entries of the identity of the preceding paragraph, obtained with the background torsion held, are not computed. On the fibre over \(y_0\) the torsion of this realisation is \((1+\epsilon)\bar T^{ab\mu}\), since \(F\to(1+\epsilon\mathsf P)F(1+\epsilon\mathsf P)^{\mathsf T}\), so the profiles of the loop terms recorded for the level-one amplitudes, with the background torsion held, are not its profiles, and its loop profiles are not computed. For the level-one realisation of the preceding paragraphs, the amplitude \(d_q\) with \(\mathsf w=\hat y\,(\hat q\cdot T)\), \(\hat y=y_0/|y_0|\) and static \(q=k\hat q\), one has \(\langle\mathsf w\cdot L(\mathsf w)\rangle=\tfrac13(\ell^2+m^2)\) and \(\sigma_1=-\tfrac13k|y_0|\epsilon\) at first order, the derivative part of \(t\) being the only one, since the frame change \(\epsilon\{\mathsf w,y^\mu\}\) is traceless; per \(\sigma_1^2\) the tree form is then \((3\rho_M/2g^2)(k^{-2}-m^2)\). In the normalisation of the identity, that is after multiplication by its constant \(L/2m^4\), which is not the operator \(L\), the leading symbol adds \(-3\varpi/m^2\) and the non-derivative term \(+3\varpi/m^4k^2\), which is not a constant at \(q=0\); the two realisations do not have the same tree form. These values are for the level-one profile \(\hat q\cdot T\). The identity of the preceding paragraph is stated at equal weights with all fibre harmonics eliminated at each fibre point; on that profile \(\langle\mathsf w\cdot L(\mathsf w)\rangle\) contains the fibre derivatives of \(L\) and is not computed. The reduced form of the paragraph on the response of the scale of \(K\) is the phase average of the induced density alone, so the term \(-3\varpi/m^2\) of the leading symbol of the tree form is not among the terms of the vacuum-term display of that paragraph or of the displays of the two preceding paragraphs.

\emph{The singlet sector at \(q=0\).} The channel is not invariant under \(L\). On the span of \(\mathsf PT\) and \(\mathsf P_\perp T\) the operator has the eigenvalue \(6\) on \(T\), the dilatation, which carries \(\sigma_1\), and \(-2\) on \((\mathsf P_\perp-3\mathsf P)T\), for which \(t=0\); on all constant symmetric \(\kappa\) it is \(-2\) on the nine traceless directions. These are non-normalisable polynomial modes, and by the localisation identity their eigenvalues are the non-derivative terms of waves with these polarisations. For \(\mathsf w=\epsilon_1T+\epsilon_y\mathsf PT\) with a common profile, \(\langle\mathsf w\cdot L(\mathsf w)\rangle=(\ell^2+m^2)(6\epsilon_1^2+2\epsilon_1\epsilon_y)\) on the fibre over \(y_0\), with \(\sigma_1=-(4\epsilon_1+\epsilon_y)\); the form is indefinite. At fixed \(\sigma_1\) it gives \(C_{\mathrm{tree}}=0\) on the channel above, \(\tfrac3{16}(\rho_M/g^2)(\ell^2+m^2)\) on the pure trace, and, at its stationary point over the \(\sigma_1\)-neutral direction, which is a maximum there, \(C_{\mathrm{tree}}=\tfrac14(\rho_M/g^2)(\ell^2+m^2)\), which in the normalisation of the identity adds \(\vartheta\varpi/2m^2>0\) to its denominator, the Yukawa sign at tree strength. This positive value is obtained by extremising over the \(\sigma_1\)-neutral direction \((\mathsf P_\perp-3\mathsf P)T\), whose own non-derivative term is negative, with the eigenvalue \(-2\) of \(L\) and the curvature \(-4\) of the form at fixed \(\sigma_1\). Which combination the static channel selects at \(q=0\) is set by the kinetic matrix of the two singlets and by the matter coupling: the leading symbol of \(\mathsf PT\) vanishes on the fibre over \(y_0\), that of \(T\) is \(\tfrac23m^2(\ell^2+m^2)k^2\) for static \(q\), and the induced and loop parts of that matrix are not computed.

\emph{Consequence.} On the two realisations of the channel along \(y^\mu y^\nu\) computed here the tree terms supply no positive constant at \(q=0\): they vanish on the rescaling, and on the level-one amplitude the leading symbol adds a negative constant and the non-derivative term is positive but of order \(k^{-2}\), whose effect on the static response at small \(k\), where the frozen base point is not controlled, is not determined. The tree non-derivative term therefore depends on the fibre profile of the wave, and its value on the profile of the identity of the preceding paragraph is not computed. A positive constant at \(q=0\) of tree strength is obtained here only on the singlet sector, at the stationary value over the \(\sigma_1\)-neutral singlet, whose own non-derivative term is negative; the identity of the preceding paragraph, written for the channel along \(y^\mu y^\nu\), does not contain it.

These statements are for the Poisson limit on \(M\), tree level, the conjugate-coordinate class with the supplied bilinear term, the matter fields at their vanishing background and \(C_{\mu\nu\lambda}=0\) on \(M\), the product background, a frozen base point with \(y^i=0\) and the parametrisation of the preceding paragraph. The localisation identity, the operator identities and the forms are exact, and the fibre moment is checked numerically. They do not compute the tree non-derivative terms on the profile of the equal-weight elimination, the loop profiles on the rescaling or on the singlet sector, the kinetic matrix of the singlets, a matter coupling or a normalisable spectrum, do not decide which realisation the static channel selects, determine no sign of the static scalar pole, and do not identify \(\sigma_1\), \(t\) or the singlets with quantities of the emergent metric of Chapter~\ref{ch:vierbein-emergence}, so no Cassini or Newtonian comparison follows. Problems~C, F and~N remain open and PROGRAMME-CONDITIONAL.

